% ============================================================
% ABSTRACT
% File: frontmatter/abstract.tex
% ============================================================

Industrial monitoring increasingly relies on data from cameras, acoustic signals and
machine sensors. Multimodal learning exploits the benefits of each data type, but simply
adding more sources does not necessarily improve prediction accuracy or reliability,
especially when data quality and relevance vary.

A review of \ReviewWorkCount{} full-text papers found
that 46 directly evaluated the value of adding an input, 21 evaluated corrupted inputs
and 11 evaluated missing inputs at inference. This thesis examines when visual
information adds value beyond machine sensors alone and whether fusion is reliable when
either stream is degraded. The initial motivation was robotic welding. An audit of the
Intel Welding data revealed that acquisition sessions were shared across partitions and
there was prior development exposure, meaning no independent welding dataset was available for
validation. Four main studies and a supporting handover study compared engineered sensor
features, pretrained visual features and patch-based temporal Transformers under
dataset-specific group splits. They evaluated concatenation, equal averaging, logistic
stacking, quality conditioned gating and visual weight transfer.

Across the 20 PHM2026
gear runs, fusing the three sensor groups to predict an unvalidated photographic proxy
resulted in a mean absolute error of 0.824 percentage points compared to 0.804 for the
training mean. In the central REASSEMBLE study (4,530 completed actions, 148
recordings), stacking video and sensor predictions increased the pooled failure average
precision  from 0.429 for sensors alone to 0.447, although the advantage over equal
averaging (0.446) was not resolved. Severe sensor noise reduced stacking AP to 0.222.
External studies did not replicate this gain. Sensor-only logistic regression
outperformed stacking in all six MSM milling protocols (0.829 vs 0.757 abnormal cutting
AP on 1,895 primary intervals), with fault-dependent quality gate behaviour. On
ViFailback (4,382 episodes, 99 held-out tasks), class-averaged AP was 0.649 for sensors
versus 0.621--0.629 for stacking across three seeds. Transferring REASSEMBLE visual
weights produced mixed results that remained below sensor level. On the 151 official
HANDOVER test trials, sensors again ranked failures above equal averaging. The value of
multimodal fusion was conditional on the validity of the target, observability of the
visual stream and sensor strength. Vision helped with failure ranking in assembly but
not in milling or manipulation under tested protocols. Adaptive weighting offered no
consistent protection against undetected corruption. These findings relate to
retrospective condition estimation and completed event classification. Early warning and
cross-factory deployment will require separate validation.
